%% ma: L10-B13-0072
%% lop: 10
%% chuong: 5
%% bai: 13
%% dang: 10.13.15
%% loai: TL
%% muc: VD
%% dap_an: "-"
%% nguon: "chuyen-de-cac-so-dac-trung-cua-mau-so-lieu-khong-ghep-nhom-toan-10-chuong-trinh-moi.pdf (GV. Phan Nhật Linh) (P03), Bài 02 Dạng 2 Bài tập tự luận 9"
%% kien_thuc: "Lớp 10 Bài 13 (số trung bình, trung vị, tứ phân vị, mốt)"
%% trang_thai: chinh-thuc
%% ghi_chu: "Ý a) lời giải gốc chọn tứ phân vị; kho chọn trung vị (số đặc trưng đo xu thế trung tâm) và giữ giá trị Q2 = 7,5 như gốc. Tự duyệt theo yêu cầu giáo viên 22/10/2026."
\begin{ex}
Hãy chọn số đặc trưng đo xu thế trung tâm phù hợp cho mỗi mẫu số liệu sau. Giải thích và tính giá trị của số đặc trưng đó.

a) Số mặt trăng đã biết của các hành tinh (theo NASA):
\begin{center}
\begin{tabular}{|c|c|c|c|c|c|c|c|c|}\hline
Hành tinh & Thủy & Kim & Trái Đất & Hỏa & Mộc & Thổ & Thiên Vương & Hải Vương\\\hline
Số mặt trăng & $0$ & $0$ & $1$ & $2$ & $63$ & $34$ & $27$ & $13$\\\hline
\end{tabular}
\end{center}
b) Số đường chuyền thành công trong một trận đấu của một số cầu thủ bóng đá: $32$; $24$; $20$; $14$; $23$.

c) Chỉ số IQ của một nhóm học sinh: $60$; $72$; $63$; $83$; $68$; $74$; $90$; $86$; $74$; $80$.

d) Các sai số trong một phép đo: $10$; $15$; $18$; $15$; $14$; $13$; $42$; $15$; $12$; $14$; $42$.
\loigiai{a) Số liệu chênh lệch rất lớn (có các giá trị $63$; $34$ lớn bất thường) nên dùng trung vị. Sắp xếp $0$; $0$; $1$; $2$; $13$; $27$; $34$; $63$, $M_e=\dfrac{2+13}2=7{,}5$.

b) Các giá trị khá gần nhau, không lặp lại nên dùng số trung bình: $\overline x=\dfrac{113}5=22{,}6$.

c) Các giá trị gần nhau, không có giá trị bất thường nên dùng số trung bình: $\overline x=\dfrac{750}{10}=75$.

d) Có giá trị bất thường ($42$) nên dùng trung vị. Sắp xếp $10$; $12$; $13$; $14$; $14$; $15$; $15$; $15$; $18$; $42$; $42$, $M_e=15$ (cũng có thể dùng mốt $15$).}
\end{ex}
